% Part: turing-machines 
% Chapter: undecidability 
% Section: unsolvability-decision-problem

\documentclass[../../../include/open-logic-section]{subfiles}

\begin{document}

\olfileid{tur}{und}{uns} 
\olsection{Masalah Putusan Ora Bisa Dirampungake}

\begin{thm}
\ollabel{thm:decision-prob}
Masalah putusan ora bisa dirampungake: Ora ana mesin Turing~$D$
kang, nalika diwiwiti ing pita kang ngemot
!!a{sentence}~$!B$ saka logika orde kapisan minangka input,
$D$~pungkasane mandheg lan ngasilake~$1$ yen lan mung yen
$!B$ sah, dene ngasilake $0$ yen ora.
\end{thm}

\begin{proof}
Upamane masalah putusan bisa dirampungake, yaiku ana
mesin Turing~$D$. Banjur kita bisa ngrampungake masalah
mandheg kaya mangkene. Kita mbangun mesin Turing~$E$ kang,
yen diwenehi input angka~$e$ saka mesin Turing~$M_e$
lan input~$w$, ngitung !!{sentence}~$!T(M_e, w) \lif !E(M_e, w)$
kang cocog, banjur mandheg karo sirah kang maca kothak
paling kiwa ing pita. Mula mesin $E \concat D$, yen diwenehi
input $e$ lan $w$, mula-mula ngitung~$!T(M_e, w) \lif
!E(M_e, w)$, banjur nglakokake mesin masalah putusan~$D$
nganggo asil iku minangka input. $D$ mandheg kanthi output~$1$
yen lan mung yen $!T(M_e, w) \lif !E(M_e, w)$ sah,
dene ngasilake~$0$ yen ora. Miturut
\olref[ver]{lem:halt-if-valid} lan
\olref[ver]{lem:valid-if-halt}, $!T(M_e, w) \lif !E(M_e, w)$
sah yen lan mung yen $M_e$ mandheg nganggo input~$w$.
Dadi, $E\concat D$, yen diwenehi input $e$ lan $w$,
mandheg kanthi output~$1$ yen lan mung yen $M_e$
mandheg nganggo input~$w$, dene mandheg kanthi
output~$0$ yen ora. Kanthi tembung liya,
$E \concat D$ bakal ngrampungake masalah mandheg.
Nanging miturut \olref[hal]{thm:halting-problem},
kita ngerti yen mesin Turing kaya mangkono ora mungkin ana.
\end{proof}

\begin{cor}\ollabel{cor:undecidable-sat}%
Ora bisa diputusake apa sembarang !!{sentence}
saka logika orde kapisan bisa dicukupi.
\end{cor}

\begin{proof}
  Upamane apa sawijining ukara bisa dicukupi iku bisa diputusake
  dening mesin Turing~$S$. Banjur kita bisa ngrampungake
  masalah putusan kaya mangkene: Yen diwenehi
  !!a{sentence}~$B$ minangka input, geser $!B$
  sakothak menyang tengen. Bali menyang kothak~$1$
  banjur tulis simbol~$\lnot$.

  Saiki lakokna mesin Turing~$S$. Mesin iki pungkasane mandheg
  kanthi output $1$ (yen $\lnot !B$ bisa dicukupi) utawa~$0$
  (yen $\lnot !B$ ora bisa dicukupi) ing pita. Yen ana
  $\TMstroke$ ing kothak~$1$, busak tandha iku;
  yen kothak~$1$ kosong, tulis $\TMstroke$, banjur mandheg.

  Mesin Turing iki mesthi mandheg, lan outpute~$1$ yen lan
  mung yen $\lnot !B$ ora bisa dicukupi, dene $0$~yen ora.
  Amarga $!B$ sah yen lan mung yen $\lnot !B$~ora bisa
  dicukupi, mesin iki ngasilake~$1$ yen lan mung yen
  $!B$ sah, dene $0$~yen ora. Tegese, mesin iki bakal
  ngrampungake masalah putusan.
\end{proof}

\begin{explain}
Dadi ora ana mesin Turing kang tansah menehi wangsulan
``ya'' utawa ``ora'' kang bener kanggo pitakon
``Apa $!B$ iku !!{sentence} logika orde kapisan kang sah?''
Nanging \emph{ana} mesin Turing kang tansah menehi
wangsulan ``ya'' kang bener---mung wae mesin iku ora
mandheg yen wangsulane ``ora''. Iki ndherek saka
teorema kasahihan lan kelengkapan logika orde kapisan,
uga saka kasunyatan yen !!{derivation}s bisa
dienumerasi kanthi efektif.
\end{explain}

\begin{thm}
  \ollabel{thm:valid-ce}%
  Kasahihan !!{sentence}s logika orde kapisan semi-bisa
  diputusake: Ana mesin Turing~$E$ kang, nalika diwiwiti
  ing pita kang ngemot !!a{sentence}~$!B$ saka logika
  orde kapisan minangka input, $E$~pungkasane mandheg lan
  ngasilake~$1$ yen lan mung yen $!B$ sah, nanging ora
  mandheg yen ora.
\end{thm}

\begin{proof}
  Kabeh !!{derivation}s logika orde kapisan kang mungkin
  bisa diasilake siji mbaka siji dening algoritma efektif.
  Mesin~$E$ nindakake iki, lan nalika nemokake
  !!a{derivation} kang nuduhake yen $\Proves
  !B$, mesin iku mandheg kanthi output~$1$.
  Miturut teorema kasahihan, yen $E$ mandheg kanthi
  output~$1$, sebabe~$\Entails !B$. Miturut teorema
  kelengkapan, yen $\Entails !B$, ana !!a{derivation}
  kang nuduhake yen~$\Proves !B$. Amarga $E$
  ngasilake kabeh !!{derivation}s kang mungkin
  kanthi sistematis, pungkasane mesin iku bakal
  nemokake kang nuduhake~$\Proves !B$,
  banjur mandheg kanthi output~$1$.
\end{proof}

\end{document}
